\documentclass[10pt,a4paper]{amsart}
\usepackage[margin=19mm]{geometry}
\usepackage{amsmath,amssymb,mathtools}
\usepackage[scr=rsfs,cal=euler]{mathalfa}
\usepackage{xfrac}
\usepackage[unicode,hidelinks,bookmarks=false]{hyperref}
\newcommand{\ie}{i.e.\ }
\newcommand{\cf}{cf.\ }
\newcommand{\resp}{respectively\ }
\newcommand{\Subsection}{\subsection}
\providecommand{\nobreakdash}{\nobreak\hbox{-}}
\setlength{\emergencystretch}{2em}
\multlinegap0pt
\usepackage{amssymb}
\usepackage{tikz}
\usepackage[normalem]{ulem}
\usepackage{bm}
\usepackage{mathtools}
\newtheorem{proposition}{Proposition}
\newcommand{\be}{\begin{equation}}
\newcommand{\cG}{\mathcal{G}}
\newcommand{\dd}{\mathrm{d}} 
\newcommand{\ee}{\end{equation}}
\title{{\bf The electromagnetic field in Poisson gauge theory: the groupoidal approach}}
\author{Fabio Di Cosmo, Vladislav G. Kupriyanov, Patrizia Vitale}
\date{Source: arXiv:2510.04858v2 (CC BY 4.0)}
\begin{document}
\maketitle

\noindent\emph{Source and licence.} {\bf The electromagnetic field in Poisson gauge theory: the groupoidal approach}. Version arXiv:2510.04858v2 (CC BY 4.0), distributed by arXiv under the Creative Commons Attribution 4.0 International licence, http://creativecommons.org/licenses/by/4.0/, as recorded in the arXiv licence metadata for this exact version and shown on the abstract page https://arxiv.org/abs/2510.04858v2. Full licence notice: \texttt{licenses/arxiv-2510.04858-CC-BY-4.0.txt}.\par\smallskip

\noindent\textit{Editorial scope.} Counted unit: the article's proposition proposition (source0.tex lines 960--966). The article's complete original proof of it is delivered with it (source0.tex lines 967--996, one \texttt{proof} environment of 30 lines). No mathematics is written, repaired or re-derived: the statement and the proof are the article's own lines, in the article's own notation, and no retained line is substituted or re-flowed.  No further source-local context is needed: every cross-reference of every retained line resolves inside the two delivered spans.  Nothing was omitted from any retained span, so the package carries no omission record.  No retained line contains a citation, so the wrapper carries no bibliography and no bibliography entry.  No retained line contains a figure environment, an image-inclusion command or drawing code, so no image is delivered and the package declares no asset dependency.  Editorial formatting only, in addition: an AMS \texttt{amsart} preamble that restates the article's own declarations and macros used by the retained lines, namely the environments \texttt{proposition} on separate counters, together with the standard packages those lines need.  The retained environments print in this document as Proposition 1 in order of appearance, whereas the article prints the same items with the numbering of its own full document.  The complete native source of the article as distributed in the arXiv e-print for this exact version is bundled unchanged as \texttt{sources/186477/source0.tex}.
\begin{proposition}
Let $\cG\rightrightarrows X$ be a local symplectic groupoid, with symplectic 2-form $\omega$, as above. 
If ${\mathcal{F}}_{jk}(y)=\left\lbrace \Pi_j,\Pi_k \right\rbrace_A\mid_{p=A}(x)$ and $F_{jk}(x)=\left\lbrace p_j-A_j,p_k-A_k \right\rbrace (x)$, we have that:
\begin{equation}
\mathcal{F}_{jk}(x) = F_{lm}(x) \left( \frac{\partial \Pi_j}{\partial p_l}\right)_{p=\mathcal{A}} \left( \frac{\partial \Pi_k}{\partial p_m} \right)_{p=\mathcal{A}}\,.    
\end{equation}
\end{proposition}

\begin{proof}
In order to prove this, let us recall that the tangent vectors to the bisection are in the kernel of the forms $\dd \Pi_{j}$ restricted to the bisection. Therefore, if $\Sigma_s$ is the embedding of $X$ given by the bisection $\Sigma=(x,\mathcal{A}(x))$, the tangent vectors 
\be
W_j = (\Sigma_s)_*\left( \frac{\partial}{\partial x^j} \right) = \frac{\partial}{\partial x^j} + \frac{\partial \mathcal{A}_k(x)}{\partial x^j}\partial_p^k
\ee
are in the kernel of the forms $\dd \Pi_j$, when restricted to the bisection $p=\mathcal{A}$. Therefore, the following conditions holds true;
\begin{equation}\label{eq.kernel}
i_{W_j}\dd \Pi_k\mid_{p=A} = \frac{\partial \Pi_k}{\partial x^{j}} \mid_{p=\mathcal{A}} + \frac{\partial \Pi_k}{\partial \mathcal{A}_{l}}\frac{\partial \mathcal{A}_l}{\partial x^{j}}\mid_{p=\mathcal{A}} + \frac{\partial \Pi_k}{\partial p_{l}}\frac{\partial \mathcal{A}_l}{\partial x^{j}}\mid_{p=\mathcal{A}} \;\;= 0
\end{equation}

Therefore, if we calculate the bracket
\begin{equation}
\left\lbrace \Pi_j, \Pi_k \right\rbrace \mid_{p=A}(x) = \gamma^{l}_{m}(x,\mathcal{A}) \left( \frac{\partial \Pi_j}{\partial x^l} \frac{\partial \Pi_k}{\partial p_m} - \frac{\partial \Pi_k}{\partial x^l}\frac{\partial \Pi_j}{\partial p_m} \right)_{p=\mathcal{A}} + \Theta^{lm}\left( \frac{\partial \Pi_j}{\partial x^l}\frac{\partial \Pi_k}{\partial x^m} \right)_{p=\mathcal{A}}\,,
\end{equation}
and we use the equation \eqref{eq.kernel} we obtain the following relation:
\begin{equation}
\begin{split}
\left\lbrace \Pi_j, \Pi_k \right\rbrace \mid_{p=A}(x) &= \left\lbrace p_l - \mathcal{A}_l, p_m-\mathcal{A}_m \right\rbrace\mid_{p=\mathcal{A}} (x)\left( \frac{\partial \Pi_j}{\partial p_l}\right)_{p=\mathcal{A}} \left( \frac{\partial \Pi_k}{\partial p_m} \right)_{p=\mathcal{A}}\, = \\ &=F_{lm}(x)\left( \frac{\partial \Pi_j}{\partial p_l}\right)_{p=\mathcal{A}} \left( \frac{\partial \Pi_k}{\partial p_m} \right)_{p=\mathcal{A}}\,.
\end{split}
\end{equation}
Finally, we have that 
\begin{equation}
\left\lbrace \Pi_j, \Pi_k \right\rbrace \mid_{p=\mathcal{A}} (x)= r_{\Sigma^{-1}}^*(\widehat{\mathcal{F}}_{jk})(x) = \mathcal{F}_{jk} (x)\,,
\end{equation}
from which one recovers that under a gauge trasformation the field strength ${\mathcal{F}}_{jk}$ transforms as 
\begin{equation}
\delta_f {\mathcal{F}}_{jk} = \mathcal{L}_{{\Theta^{\sharp}(\dd f)}} {\mathcal{F}}_{jk} = \left\lbrace f, {\mathcal{F}}_{jk} \right\rbrace_{\Theta}\,,
\end{equation}
because $\widehat{\mathcal{F}}_{jk}(y)$ are gauge-invariant objects. This shows that $\mathcal{F}_{jk}(x)$ has the right transformation properties under gauge transformations.
\end{proof}

\end{document}
